\documentclass[a4paper]{article}

\usepackage[english]{babel}
%\usepackage[utf8x]{inputenc}
\usepackage{amsmath}
\usepackage{graphicx}
\usepackage[colorinlistoftodos]{todonotes}
\usepackage[all]{xy}
\usepackage{dsfont}
\usepackage{amssymb}
\usepackage{bbold}
\usepackage{amsthm}

\usepackage{tikz}
\usepackage{verbatim}
\usetikzlibrary{matrix}
\usetikzlibrary{cd}

\usepackage[autostyle=true]{csquotes}
\usepackage{pdfpages}
\usepackage{hyperref}
\usepackage{fancyhdr}
\usepackage{multicol}



\newtheorem{theorem}{Theorem}[section]
\newtheorem{proposition}{Proposition}[theorem]
\newtheorem{lemma}{Lemma}[theorem]
\newtheorem{corollary}{Corollary}[theorem]
\newtheorem*{definition}{Definition}


\newcommand{\norm}[1]{\left\lVert#1\right\rVert}

\newcommand{\Rn}{\mathbb{R}^n}


\newcommand{\too}[1]{\overset{#1}{\longrightarrow}}
\newcommand{\hookar}[1]{\overset{#1}{\hookrightarrow}}
\newcommand{\colim}[1]{ \underset{#1}{\texttt{colim}} \text{ } }
\newcommand{\fset}[1]{ \{ #1_0 , \dots , #1_n \} }
\newcommand{\ddt}{\frac{d}{dt}}


\newcommand{\N}{\mathbb{N}}
\newcommand{\Z}{\mathbb{Z}}







\fancyhead[L]{\includegraphics[scale=0.1]{UFT.png}}

\begin{document}


\thispagestyle{fancy}
\vspace*{1cm}

\begin{minipage}[r]{12cm}
    \begin{flushright}
\begin{tikzpicture}
\draw[green, very thick] (0,0) rectangle (3,2);
\end{tikzpicture}
    \end{flushright}
\end{minipage}

\begin{center}
\Large{\textbf{Control 1\\ Cálculo Vectorial\\ }}
\end{center}
\vspace{1cm}
\begin{minipage}[r]{15cm}
\begin{flushleft}

Fecha: 10 de septiembre 2025 \\%
Duración: 120 minutos  \\

\end{flushleft}
\end{minipage}
\vspace{2cm}
\begin{center}
\Large{\textbf{Facultad de Ingeniería}}
\end{center}
\vspace{1cm}
\setlength{\parindent}{0cm}
“Es deber de cada alumno mantener una actitud moral y éticamente intachable en el transcurso de toda evaluación” \\

En caso de que un alumno sea sorprendido en actitudes deshonestas ya sea intentando o
dejándose copiar, será calificado con la nota mínima (1.0) y perderá toda posibilidad de
eliminar dicha nota. Esta conducta quedará archivada como antecedente en su carpeta y podrá
constituir causal de eliminación. Asimismo, está \textbf{prohibido salir de la sala, usar o tener
teléfonos celulares (en bolsillos o sobre las mesas)} o cualquier tipo de aparato electrónico,
tales como audífonos, smartwatch, entre otros, a no ser que el profesor lo autorice.\\


Escala al $60\%$ \\

Puntaje para un 4.0: 36 puntos


\vspace{3cm}
\makebox[12cm]{\textbf{Nombre Alumno}:\ \hrulefill} 
\newpage
\thispagestyle{empty}

\setlength{\parindent}{0cm}

\textbf{Ejercicio 1 \textbf{(12 puntos)}:} Responda las siguientes preguntas \\


a) Determine la ecuación general del plano que pasa por el punto $(1,1,1)$ y que contiene a la recta de ecuación vectorial $L_1: (2,0,1)+t(1,-1,2)$. \\

\ 

c) Demuestre que el plano  encontrado en a) es paralelo a la recta $L_1$ de ecuación general 

\begin{align*}
    \frac{x-1}{3} = \frac{1-y}{3} =\frac{2z+3}{5}
\end{align*}


\newpage


\phantom{xc}


\newpage

\textbf{Ejercicio 2 \textbf{(12 puntos)}:} Calcule los siguientes límites, en caso de no existir demuéstrelo \\

a) $\displaystyle \lim_{(x,y) \to (0,0)} \frac{x^2y}{\sqrt{x^2+y^2}}$ \\


b) $\displaystyle \lim_{(x,y) \to (3,3)} \frac{9x^2-y^4}{3x-y^2}$ \\

c) $\displaystyle \lim_{(x,y) \to (0,0)} \frac{x^3-3x^2y}{(x^2+y^2)^{3/2}}$ \\ \\ \\

\newpage


\phantom{xc}

\newpage


\textbf{Ejercicio 3 \textbf{(12 puntos)}:} Considere el cilindro $x^2+y^2=16$ y el plano $2x-3y+z=4$, encuentre una parametrización $r(t)$ de la intersección de estas dos superfices, luego demuestre que $r(0)$ es perpendicular a la recta $L$ de ecuación

\begin{align*}
    \frac{x-1}{2} = \frac{1-y}{5} = \frac{3z-2}{6}
\end{align*}

\newpage


\phantom{xc}


\newpage


\textbf{Ejercicio 4 \textbf{(12 puntos)}:} Considere la función dada por 

\begin{align*}
    f(x,y) = \ln(x^2-y)
\end{align*}

a) Determine el dominio de $f$ y luego grafíquelo. \\

b) Determine la curva de nivel $C_0(f)$ y grafíquela.


\newpage


\phantom{xc}

\newpage

\textbf{Ejercicio 5 \textbf{(12 puntos)}:} Sea $\alpha \in \mathbb{R}$, considere la función

\begin{align*}
    u(x,y) = e^{-\alpha y} \sin(2x)
\end{align*}

Demuestre que $u$ satisface 

\begin{align*}
    u_y-\alpha u_{xx} =0
\end{align*}




\newpage

\phantom{xc}

\newpage

\begin{center}
\textbf{Formulario}
\end{center}

\bigskip

\textbf{1. Longitud entre dos vectores $v$ y $w$} \\[1mm]


 
\[
d(v,w) = || v-w||
\]

Ángulo entre dos vectores


\[
\theta = \arccos\left( \frac{v \cdot w}{||v|| ||w||}\right)
\]

- Ecuación vectorial de la recta con vector posicicón $p$ y vector director $d$ 
\[
L : p+tv, \text{ con } t\in \mathbb{R}
\]

\bigskip

\textbf{2. Ecuación general del plano} \\[1mm]


\[
Ax+By+Cz+D=0
\]

con $n=(A,B,C)$. \\




\textbf{Coordenadas Polares}

\begin{align*}
    x&= r\cos(\theta) \\
    y&=r\sin(\theta)
\end{align*}



\end{document}